\scanpage{201}

% B01 / text / top, structure (14)
\mspage{103}$\cdot$\quad Pour que la structure (14)

% B02 / display / top, (17), glosses on right
\[
 (S,S^*,R_*,\varphi_0)\quad\left|\begin{aligned}
 &S\ \text{``ens des sommets''}\\
 &S^*\ \text{ens des ``étoiles'' ou ``cosommets''}\\
 &R_*\ \text{relation ``d'incidence''}\\
 &\varphi_0\ \text{\deleted{application} }\underset{\inserted{\text{jauge}}}{\text{application}}
 \end{aligned}\right. \tag*{(17)}
\]

% B03 / text / upper, conditions for realization
provienne, à isom près, d'une structure de graphe \inserted{$A_0\subset\unsure{P|\mathcal P|\mathfrak P}_2(S)$}
sur $S$, \deleted{\unsure{ou d'un|i.e. d'une}} satisfaisant (13), il faut évidemment
qu'elle satisfasse les conditions suivantes

% B04 / display / upper middle, (18), brace covering a and b
\[
 \left\{\begin{aligned}
 &\text{a) Les applications}\\
 &\qquad s^*\longmapsto{}^sR_*(s^*)\qquad S^*\longrightarrow P(S)\\
 &\qquad s\longmapsto R_*(s)\qquad S\longrightarrow P(S^*)\\
 &\qquad\text{sont \uline{injectives}}\\
 &\text{b) L'application }\varphi_0:S\longrightarrow S^*\text{ est \uline{bijective}, et}\\
 &\quad\text{satisfait la condition de symétrie}
 \end{aligned}\right. \tag*{(18)}
\]
% slanted
\margin{\uline{NB} Si on n'exige pas que $(S,A_0)$ satisfasse la condition (13), il faut laisser tomber dans a) l'injectivité de $S\longrightarrow P(S^*)$, et dans b) l'injectivité de $\varphi_0$ -- $\varphi_0$ sera seulement surjective.}

% B05 / text / middle, symmetry and NB
La relation $(s,\varphi_0(t))\in R_*$ en $(s,t)$ est symétrique
i.e. $\foralls,t\in S$, on a \deleted{$s\ne t$} l'équivalence
\deleted{\unsure{(19)|(18)}}\margin{\uline{NB} On peut remplacer la donnée de $\varphi_0$, plus symétriquement, par la donnée d'une 2\textsuperscript{ième} relation $\Phi_0\subset S\times S^*$, bijective $\ldots$}

% B06 / display / middle, (19)
\[
 (s,\varphi_0(t))\in R_*\Longleftrightarrow(t,\varphi_0(s))\in R_* \tag*{(19)}
\]

% B07 / text / middle, c), and parenthesis
c) \deleted{La relation envisagée \unsure{en}}\inserted{\unsure{est sans}} $\forall s\in S$, on a $(s,\varphi_0(s))\notin R$.
(L'injectivité de $\varphi_0:S\longrightarrow S^*$ i.e. de l'application composée
$S\longrightarrow S^*\longrightarrow P(S)$ n'est autre que l'hypothèse d'injectivité (13).

% B08 / text / middle below, reciprocal
La réciproque est évidente. Ainsi, on trouve
\deleted{Nous avons disti}

% B09 / diagram / lower, (20), groupoid equivalence
\diagram{(20)\quad Équivalence de groupoïdes\par
\begin{tikzcd}[ampersand replacement=\&,column sep=large]
 {\substack{\text{(A)}\\\text{catégorie-groupoïde}\\\text{des graphes stricts}\\(S,A_0)\text{ satisfaisant (13)}\\\text{(sommets déterminés par}\\\text{leurs étoiles)}}}\arrow[r,"\approx"] \&
 {\substack{\text{(B)}\\\text{catégorie-groupoïdes des}\\\text{systèmes }\deleted{\unsure{(S,R)}}(S,S^*,R_*)\\\text{``sommets-cosommets-incidence''}\\\text{\deleted{satisfaisant} condition}\\\text{d'inj. symétrique, munie}\\\text{d'une }\varphi_0\,\text{application bijection}\\S\xrightarrow{\varphi_0}S^*\quad\text{``symétrique''}\\\text{\deleted{et} ``antiréflexive'' }((s,\varphi_0(s))\notin R_0)}}
\end{tikzcd}}

% B10 / text / lower, duality and left-hand gloss
La description (B) fait apparaître une \uline{dualité} dans
la catégorie des graphes envisagés, (``graphes stricts séparés\deleted{\unsure{ants}}''), via
% slanted
\margin{``séparé'' signifie que la notion d'étoile d'un sommet permet de ``séparer'' les sommets''}

% B11 / display / lower, (21), terms renamed below
\[
 (S,S^*,R_*,\varphi_0)\longmapsto
 (\underset{\substack{\shortparallel\\S'}}{S^*},
 \underset{\substack{\shortparallel\\S'{}^*}}{S},
 \underset{\substack{\shortparallel\\R'_*}}{{}^sR_*},
 \underset{\substack{\shortparallel\\\varphi'_0}}{\varphi_0^{-1}}), \tag*{(21)}
\]

% B12 / text / foot, dual graph
d'où la notion de ``graphe dual'' d'un graphe
strict séparé ; mais celui-ci est \uline{canoniquement} isomorphe
au graphe de départ, par

% B13 / display / foot, (22) and right marginal note
\[
 S\xrightarrow{\varphi_0}_{\sim} S^*\qquad S^*\xrightarrow{\varphi_0^{-1}}_{\sim}S \tag*{(22)}
\]
\margin{\uline{NB} La condition de \uline{symétrie} (18) b) signifie justement que (22) est compatible avec les relations $\ldots$}

% archive sheet number: 200
